Для решения методом Ньютона исходная система приводится к стандартному виду
\[
    F(X)=0,
\]
где
\[
    X=
    \begin{pmatrix}
        x \\
        y
    \end{pmatrix},
    \qquad
    F(X)=
    \begin{pmatrix}
        f_1(x,y) \\
        f_2(x,y)
    \end{pmatrix}.
\]
Для рассматриваемой задачи вводятся функции
\[
    f_1(x,y)=\cos(y-1)+x-0.5,
    \qquad
    f_2(x,y)=y-\cos(x)-3.
\]
Тогда решение исходной системы эквивалентно решению системы
\[
    \begin{cases}
        f_1(x,y)=0, \\
        f_2(x,y)=0.
    \end{cases}
\]

Метод Ньютона позволяет построить последовательность приближений
\[
    X_0,\;X_1,\;X_2,\;\dots,
\]
сходящуюся к решению системы при удачном выборе начального приближения. Итерационная формула метода имеет вид
\[
    X_{k+1}=X_k-\bigl(F'(X_k)\bigr)^{-1}F(X_k),
    \qquad k=0,1,2,\dots
\]
Здесь \(F'(X_k)\) $-$ матрица Якоби, составленная из частных производных функций \(f_1\) и \(f_2\).

Для данной системы матрица Якоби равна
\[
    J(x,y)=
    \begin{pmatrix}
        \frac{\partial f_1}{\partial x} & \frac{\partial f_1}{\partial y} \\
        \frac{\partial f_2}{\partial x} & \frac{\partial f_2}{\partial y}
    \end{pmatrix}
    =
    \begin{pmatrix}
        1       & -\sin(y-1) \\
        \sin(x) & 1
    \end{pmatrix},
\]
так как
\[
    \frac{\partial f_1}{\partial x}=1,
    \qquad
    \frac{\partial f_1}{\partial y}=-\sin(y-1),
    \qquad
    \frac{\partial f_2}{\partial x}=\sin(x),
    \qquad
    \frac{\partial f_2}{\partial y}=1.
\]

Обратная матрица Якоби явно не вычисляется. Вместо этого на каждой итерации решается система линейных алгебраических уравнений относительно поправок
\[
    \Delta X_k=
    \begin{pmatrix}
        \Delta x_k \\
        \Delta y_k
    \end{pmatrix}:
\]
\[
    J(X_k)\,\Delta X_k=-F(X_k).
\]
После нахождения поправок новое приближение определяется по формулам
\[
    x_{k+1}=x_k+\Delta x_k,
    \qquad
    y_{k+1}=y_k+\Delta y_k.
\]


Пересечения кривых, найденная графически, имеет координаты
\[
(x_0,\;y_0)\approx(1.20691,\;3.35591).
\]

В качестве начального приближения использовано близкое к решению на графике значение
\[
    X_0=
    \begin{pmatrix}
        1 \\
        3
    \end{pmatrix}.
\]

Критерием окончания итераций служит достижение заданной точности. В работе используется норма вектора поправок
\[
    \|\Delta X_k\|=\max\left(|\Delta x_k|,\;|\Delta y_k|\right).
\]
Если выполнено условие
\[
    \max\left(|\Delta x_k|,\;|\Delta y_k|\right)\leq \varepsilon,
\]
то итерационный процесс прекращается.